\documentclass{report}
\usepackage{amssymb}
\begin{document}
\title{libdmagnetic2- API description}
\author{Thomas Dettbarn}
\maketitle
\tableofcontents
\newpage
\chapter{Introduction}
The purpose of this chapter is to introduce the reader to conecpts of the library.
\section{Mode of operation}
EXPLAIN: There are a total of three libraries. One for loading/converting input sources. One for running the engine. One for decoding the pictures/animations. 
Each library has a function \verb+get_size()+, which returns the number of bytes required for the handle. 
\begin{verbatim}
void *pHandle;
int errcode;
errcode=dMagnetic2_engine_get_size(&bytes);
pHandle=malloc(bytes);
errcode=dMagnetic2_engine_init(pHandle);
\end{verbatim}
\section{Helloworld example}
%%%%%%%%%%%%%%%%%%%%%%%%%%%%%%%%%%%%%%%%%%%%%%%%%%%%%%%%%%%%%%%%%%%%%%%%%%
\chapter{libdmagnetic2 loader.a}
The purpose of this chapter is to describe how to use the API to load Magnetic Scrolls Games from various sources.
\section{Helloworld example}
The purpose of this section is to provide a source code of a helloworld example which converts input sources into the .mag/.gfx format.
For now, see \verb+tests/loader/loader_mkmaggfx.c+

%%%%%%%%%%%%%%%%%%%%%%%%%%%%%%%%%%%%%%%%%%%%%%%%%%%%%%%%%%%%%%%%%%%%%%%%%%
\chapter{libdmagnetic2 engine.a}
The purpose of this chapter is to describe how to use the API to run Magnetic Scrolls Games.
\section{Helloworld example}
The purpose of this section is to provide a source code of a helloworld example which runs a simple (playable) engine.
For now, see \verb+tests/engine/engine_runmag.c+.
\section{Data structures}
\section{Interface functions}
\section{Saving and Loading games}
Get the context from the handle by calling \verb+dMagnetic2_engine_save_game();+. Overwrite the context in the handle by calling \verb+dMagnetic2_engine_load_game();+.

%%%%%%%%%%%%%%%%%%%%%%%%%%%%%%%%%%%%%%%%%%%%%%%%%%%%%%%%%%%%%%%%%%%%%%%%%%
\chapter{libdmagnetic2 graphics.a}
The purpose of this chapter is to descripe how to use the API to decode the graphics and animations.
\section{Helloworld example}
The purpose of this chapter is to provide a source code for decoding graphics and storing them as .xpm files.
For now, see \verb+tests/graphics/graphics_animations.c+, \verb+tests/graphics/graphics_pictures1.c+, \verb+tests/graphics/graphics_pictures2.c+.
\section{Data structures}
\section{Interface functions}
\end{document}
